\begin{abstract}
Large language model (LLM) services are deployed on open-source serving engines whose
schedulers decide, under load, which requests run, which wait, and which are evicted.
These decisions are invisible to the application, yet they determine what users experience.
Published comparisons of serving engines mostly report peak throughput or per-token speed,
which says little about how an engine behaves once its memory or scheduling capacity is
exhausted. We present an empirical comparison of three widely used engines, \vllm, \sglang,
and \trt, driven through their own OpenAI-compatible frontends by a common open-loop
benchmark harness. We measure six user-facing metrics (time to first token, time between
tokens, time per output token, end-to-end latency, throughput, and request outcome) and
combine them into SLO goodput, across three models with different key--value (KV) cache
footprints, realistic mixed and multi-turn workloads, five arrival processes, two scheduler
capacity settings, two prefill-handling features, n-gram speculative decoding, and two
hardware platforms.

The engines are nearly indistinguishable at light load and diverge sharply once a resource
saturates, and they diverge in \emph{where} they make requests wait. On heterogeneous
traffic the engine with the fastest token generation delivers the lowest goodput, while on
uniform traffic it is the fastest end to end; a performance feature that helps one engine
can hurt another. We also report methodology pitfalls that change conclusions: time to first
token measured from dispatch hides client-side queueing, finite traces make latency
non-monotonic in load, and identically named engine metrics are not comparable. We conclude
with a decision guide for practitioners and release the harness and data.
\todo{tighten to the final RQ findings; ICPE abstracts are typically 150--250 words}
\end{abstract}
